\section{An Entropy Tangent for the Angular Error}
\label{app:centered-support}

This appendix proves \eqref{eq:centered-support}. Its purpose is
to preserve equality at the dictator's two posterior values
\(\pm r\), while controlling the error of the line
\(x\mapsto(\arcsin r)x/r\).

Fix \(5/4\le\lambda\le\pi/2\). On \(0\le x\le1\), define
\begin{align*}
 A_{\rm sc}(x)&=x\arcsin x,\\
 e(x)&=\arcsin x-\lambda x,\\
 \Psi(x)&=\phi(x)+e(x)^2/2.
\end{align*}
We will show that \(\Psi\) increases and that \(A_{\rm sc}\) is
convex as a function of \(\Psi\). A tangent at \(x=r\) will then
give exactly the required inequality.

Put \(x=\sin t\), \(c=\cos t\), and \(g=\atanh(\sin t)\).
Dots in this paragraph denote derivatives in \(t\). Direct
differentiation gives
\begin{align*}
 \dot A_{\rm sc}&=ct+x,& \ddot A_{\rm sc}&=2c-xt,\\
 \dot\Psi&=gc+e(1-\lambda c),\\
 \ddot\Psi&=1-gx+(1-\lambda c)^2+\lambda xe.
\end{align*}
The terms linear in \(\lambda\) cancel in the curvature numerator:
\[
 K_\lambda:=\ddot A_{\rm sc}\dot\Psi-
                \dot A_{\rm sc}\ddot\Psi
 =g(1+c^2)-x(t^2+2)+\lambda^2(x-c^3t).
\]
Since \(t<\tan t=x/c\), its last coefficient is positive.
It suffices to show \(K_{5/4}>0\).

Here are two matched integral bounds, with their remainders exposed:
\begin{align*}
 \left[\frac{y(15+4y^2)}{15+9y^2}-\arctan y\right]'
   &=\frac{36y^6}{(15+9y^2)^2(1+y^2)},\\
 \left[\atanh y-\frac{y(15-4y^2)}{15-9y^2}\right]'
   &=\frac{36y^6}{(15-9y^2)^2(1-y^2)}.
\end{align*}
Both differences vanish at zero. Substituting \(y=\tan(t/2)\)
and using the half-angle formulas gives
\[
 \frac tx\le\frac{19+11c}{3(1+c)(4+c)},\qquad
 \frac gx\ge\frac{11+19c}{3(1+c)(1+4c)}.
\]
Use these upper bounds for the negative \(t,t^2\) terms and the
lower bound for \(g\). After multiplication by positive
denominators the result is
\[
 \frac{16K_{5/4}}x\ge
 \frac{(1-c)Q(c)}{9(1+c)(4+c)^2(1+4c)},
\]
where the following identity proves the sign directly:
\begin{align*}
 Q(c)={}&375(2-6c-7c^2)^2+164(1-c)^5\\
 &+740c(1-c)^4+c^2(1-c)^3\\
 &+25c^3(1-c)^2+5925c^4(1-c)+8625c^5.
\end{align*}
Indeed expansion gives
\(Q(c)=1664-9080c+1681c^2+34322c^3+22113c^4+3300c^5\),
the numerator produced by the preceding substitution.
Each term in the identity is nonnegative on \(0\le c\le1\),
and not all vanish. Thus \(K_\lambda>0\) in the interior.

To justify using \(\Psi\) as a coordinate, observe that
\[
 \left(\frac{\dot\Psi}{\dot A_{\rm sc}}\right)'
       =-\frac{K_\lambda}{\dot A_{\rm sc}^{\,2}}<0.
\]
As \(t\uparrow\pi/2\), this ratio tends to \(\pi/2-\lambda\ge0\).
It is therefore positive before the endpoint. This proves
\(\dot\Psi>0\), and
\(d^2A_{\rm sc}/d\Psi^2=K_\lambda/\dot\Psi^3>0\).

Now let \(r=\tanh\ell\), \(\ell\ge31/20\), and
\(\lambda=\arcsin(r)/r\). This slope lies in \([5/4,\pi/2]\).
For the lower bound, \(e^{31/10}>20(1+1/10)=22\) gives
\(r>21/23\), hence \(r^2>5/6\).
The alternating sine bound, with \(y=5r/4<5/4\), gives
\[
 \frac{\sin(5r/4)}r
 \le\frac54-\frac{125}{384}r^2+\frac{3125}{122880}r^4<1.
\]
The quadratic on the right decreases on \(0\le r^2\le1\);
its value at \(5/6\) is below one. As \(5r/4<\pi/2\),
\(\arcsin r>5r/4\). The upper bound is the endpoint chord
of the convex function \(\arcsin r\).

At \(x=r\) we have \(e(r)=0\), \(\Psi(r)=\phi(r)\), and
\[
 \frac{dA_{\rm sc}}{d\Psi}(r)
   =\frac{\arcsin r+r/\sqrt{1-r^2}}{\atanh r}
   =\alpha.
\]
The supporting line of the proved convex function is
\(A_{\rm sc}(x)-A_{\rm sc}(r)\ge
\alpha[\phi(x)-\phi(r)+e(x)^2/2]\).
This is \eqref{eq:centered-support}; evenness extends it to
negative \(x\), and continuity includes the endpoints.
